\section{Background}
\label{sec:background}

\subsection{Attention variants}
For a sequence of length $N$ with head dimension $d$, self-attention computes
$O = \mathrm{softmax}(QK^\top/\sqrt{d})V$. During decode, the query is a single
new token ($N_q{=}1$) attending over a KV cache of length $N_k$. The variants
differ in how much KV must be read per generated token:
\begin{itemize}
  \item \textbf{MHA} (multi-head): $H_q{=}H_{kv}$; every query head has its own KV,
  so the cache is $O(B \cdot H \cdot N_k \cdot d)$ and decode arithmetic intensity
  (\ai, FLOPs per byte) is $\approx 2/b$ where $b$ is bytes per KV element.
  \item \textbf{GQA} (grouped-query~\cite{gqa}): $H_{kv} < H_q$, with
  $G{=}H_q/H_{kv}$ query heads sharing one KV group. A decode kernel can pack the
  $G$ queries of a group into one CTA, reading the shared KV \emph{once} and
  raising \ai\ from $2/b$ to $2G/b$.
  \item \textbf{MLA} (multi-head latent, DeepSeek~\cite{deepseekv2,deepseekv3}):
  one shared low-rank latent (dimension $576 = 512$ latent $+\,64$ RoPE) replaces
  all KV heads. At decode all $h_q$ query heads attend to the single latent, so
  the packed row dimension is $M{=}h_q{=}128$ \emph{by construction}, and
  $\ai \approx 2\,h_q(2L{+}R)/((L{+}R)\,b) \approx 3.78\,h_q/b$. MLA is the only
  decode shape that meets the Blackwell block-scaled NVFP4 tensor-core gate
  ($M \ge 128$) with no speculative draft.
\end{itemize}

\subsection{The roofline model, and why it fails for decode}
The classical roofline~\cite{roofline} bounds attainable throughput by
$\min(\text{peak FLOP/s},\ \ai \times \text{peak bandwidth})$; the ridge point
$\ai^\star = \text{peak FLOP/s} / \text{peak bandwidth}$ separates memory-bound
($\ai < \ai^\star$) from compute-bound ($\ai > \ai^\star$). Applied naively to
modern GPUs, this over-predicts sustained throughput by wide margins: datasheet
peaks are unreachable, and the model is blind to the L2 cache, per-CTA
scheduling, occupancy, and tensor-core pipeline fill. Microbenchmark-driven
analysis~\cite{roofmethod} quantifies this: naive roofline exceeds 95\% error on
B200 GEMMs where a microbenchmark-grounded model reaches 1.31\% MAE. Our
per-CTA-corrected layer (Section~\ref{sec:method}) adds exactly those on-chip
constraints: shared-memory staging $\to$ occupancy, warp packing (GEMV vs GEMM),
and pipeline depth vs available MMA tiles.

\subsection{Blackwell Ultra (B300 / \texttt{sm\_103})}
\label{sec:bg-blackwell}
Our headline architecture is the B300 (GB300, \texttt{sm\_103}). We use
\emph{measured} constants throughout, several of which correct commonly cited
datasheet values (Table~\ref{tab:archs}):
148 streaming multiprocessors (SMs) enabled (not 160), 8\tbs\ HBM (flat versus
B200; only capacity grew to 288\,GB), \textbf{132.6\,MB} L2 (measured; the widely
repeated ``192\,MB'' appears to conflate B200's 192\,GB capacity), 228\,KB
shared memory per SM, and a measured boost clock of \textbf{2032\,MHz} (below the
2600\,MHz marketing figure). The \texttt{tcgen05} tensor cores deliver dense
peaks of 2.5\,PF (FP16), 5\,PF (FP8), and 15\,PF (block-scaled NVFP4). We also
measured the exponential throughput used by softmax: an EX2 microkernel sustains
\textbf{5.33\,TExp/s}, only $0.50\times$ the 10.7\,TExp/s vendor peak. We fold
this datasheet gap into the model, though it is immaterial to $M{=}1$ decode
(the MUFU term is $<3\%$ of decode time).

\begin{table}[t]
  \centering
  \caption{Measured architectural constants (from \texttt{roofline/archs.py}).
  Blackwell values are torch device-property / microbenchmark measurements taken
  on the rented parts; peaks in PF are dense tensor-core datasheet figures.}
  \label{tab:archs}
  \footnotesize
  \setlength{\tabcolsep}{3.5pt}
  \begin{tabular}{lrrrr}
    \toprule
    & T4 & A100 & B200 & B300 \\
    & \texttt{sm\_75} & \texttt{sm\_80} & \texttt{sm\_100} & \texttt{sm\_103} \\
    \midrule
    SMs                 & 40    & 108   & 148   & 148 \\
    HBM BW              & 320\gbs & 2.0\tbs & 8\tbs & 8\tbs \\
    HBM capacity        & 16\,GB & 80\,GB & 192\,GB & 288\,GB \\
    L2 cache            & 4\,MB & 40\,MB & 132.6\,MB & 132.6\,MB \\
    Boost clock         & 1590 & 1410 & 1965 & 2032\,MHz \\
    \smem/SM            & 64\,KB & 164\,KB & 228\,KB & 228\,KB \\
    FP16 TC peak        & 65\,TF & 312\,TF & 2.25\,PF & 2.5\,PF \\
    FP8 TC peak         & --    & --    & 4.5\,PF & 5\,PF \\
    NVFP4 TC peak       & --    & --    & 9\,PF & 15\,PF \\
    \bottomrule
  \end{tabular}
\end{table}
